\section{Related Work}
\label{sec:related}

\paragraph{Automated harness optimization.}
The closest line of work uses LLMs as outer-loop optimizers for prompts,
agent programs, or harness code. OPRO~\cite{yang2024opro} casts prompt
engineering as black-box optimization with LLM-generated candidates;
GEPA~\cite{agrawal2025gepa} refines prompts using reflective textual
feedback. ADAS~\cite{hu2024adas} moves from prompts to agent programs,
having an LLM propose new scaffolds from an archive of prior designs.
Meta-Harness~\cite{lee2026metaharness} is closest to our setting: it
optimizes harness files with a coding agent that can inspect the full
filesystem history of previous candidates. Related program-search systems
such as FunSearch, AlphaEvolve, and OpenEvolve use LLMs as mutators under
automated fitness functions; we discuss these broader connections in
Appendix~\ref{app:supp_related}. \name keeps this file-backed optimization
interface but changes the information policy around it. Rather than exposing
the entire accumulated workspace or compressing it into a fixed summary,
\name chooses which prior iterations and file hints should be salient before
each proposal.

\paragraph{Context curation in iterative optimization.}
Our motivation aligns with evidence that more context is not always better.
Liu et al.~\cite{liu2024lost} show that models use long contexts unevenly;
CODEFILTER~\cite{li2025codefilter} demonstrates that filtering retrieved code
context can outperform passing all retrieved chunks. Nie et
al.~\cite{nie2026learningloop} further show that iterative generative
optimization is sensitive to design choices such as credit horizon and
experience batching. \name carries these ideas into automated harness
optimization by treating proposer attention as an online control variable: a
progressive iteration policy manages history scope, and delayed utility
estimates manage file emphasis.
